% Lösungen Einheit 3 von 6 – Einheit 3 (zusammenbau v0.1)
\einheitenkopf{Einheit 3 von 6 · Einheit 3}

\erg{15}{a) $D = ]3; \infty[$, denn $x - 3 > 0$ \quad b) $W = ]2; \infty[$ \quad c) Schranke $5$: $2e^{-x}$ ist für jedes $x$ positiv, von $5$ wird also immer etwas abgezogen \quad d) $W = ]0; 1]$, denn $-x^2 \le 0$ \quad e) $W = [1; 3]$; der Wert $2$ zweimal (bei $0$ und bei $\pi$) \quad f) $D = \mathbb{R}$, da $kx^4 + 1 \ge 1 > 0$; $g_k(0) = \mathrm{ln}(1) = 0$ für jedes $k$; $k + 1 = e^3$, also $k = e^3 - 1 \approx 19{,}09$}
\erg{16}{a) Fehler: der Wert $4$ wird nie angenommen, weil $e^{-x} > 0$. Richtig: $W = ]-\infty; 4[$ \quad b) $\mathrm{ln}(x)$ ist die Hochzahl, mit der man $e$ potenzieren muss, um $x$ zu erhalten; jede Potenz von $e$ ist positiv, für $x \le 0$ gibt es keine solche Hochzahl. \quad c) ab $t = 2\,\mathrm{ln}(16) \approx 5{,}55$ h ist $T \le 5$ °C; $4$ °C werden nie erreicht, weil $16e^{-0{,}5t} > 0$}
